\begin{textsample}{NEG Dim 2 – al – Score -1.00 – al/al\_ef\_33.txt}  \label{ex:f2_neg_031}
4.
METODOLOGIAS E ABORDAGENS NO COMPONENTE CURRICULAR ARTE As metodologias e abordagens estimulam percepções e encontros diferenciados ao valorizar habilidades, competências e experiências específicas.
Analisando o que diferencia metodologia de abordagem, percebe -se que as metodologias são diretrizes que orientam os processos de ensino e aprendizagem e que se concretizam em estratégias, abordagens e técnicas concretas, específicas, diferenciadas, ou seja, metodologia é como se estabelece a relação entre sujeitos e objetos de ensino e pesquisa.
Enquanto que a abordagem tem como base a transculturalidade, que valoriza a identidade dos diversos povos, sociedades e práticas culturais, concentrando -se nos diálogos, encontros e construções conjuntas das diversas culturas e tradições, valorizando o surgimento do novo e de novas identidades culturais, ou seja, é como chegar nos sujeitos e suas culturas híbridas.
É importante considerar articulações metodológicas que mobilizem ações e experiências significativas que estimulem o interesse do estudante em aprender, de forma que o processo seja conduzido em direção ao conhecimento, ampliando o que se conhece, conhecendo o desconhecido, conquistando a autonomia nos modos de perceber e interferindo na realidade.
A Abordagem \textbf{Triangular}, nos termos de Ana Mae Barbosa ( 2010 ), considerada a base da maioria dos programas de Arte-Educação, possui três eixos de aprendizagem, sem ordem preestabelecida : apreciar, contextualizar e fazer.
Tais eixos podem ampliar a capacidade cognitiva e crítica dos sujeitos aprendentes, estimulando -os a criarem suas próprias manifestações poéticas e artísticas com um repertório cultural alimentado pelas próprias produções.
Hoje esses eixos passaram a ser chamados dimensões e conforme a BNCC ( 2017 ), ampliadas com mais três, formando assim as seis dimensões do conhecimento : criação, crítica, estesia, expressão, fruição e reflexão.
Outras concepções de educação estética, artística e cultural, no entanto, vêm trilhando caminhos próprios nas escolas e nos programas educativos.
A interdisciplinaridade proposta por Fernando Hernández com a ideia de cultura visual busca referenciais de arte, arquitetura, história, mediação cultural, psicologia e antropologia, que não se organiza somente com base em nomes de peças, fatos e sujeitos, mas na relação estabelecida com seus significados culturais.
Defende -se, de acordo com Hernández ( 2000, p. 54 ), uma abordagem que considere “ a arte e a cultura como mediadores de significados, na qual o significado pode ser interpretado e construído, as imagens podem informar àqueles que as veem sobre eles mesmos e sobre temas relevantes no mundo ”.
As aprendizagens artísticas exploram diversos conceitos que visam potencializar a experiência com arte.
Conceitos de forma e conteúdos que articulados aos elementos de linguagens criam estilos, discursos e poéticas.
Esse processo criador, nos termos de Vygotsky ( 1979 ), ao interpor a realidade, a imaginação, a emoção e a cognição, envolve reconstrução, reelaboração e redescoberta.
Ensinar Arte significa proporcionar aos estudantes o conhecimento a partir de um modo formativo, inventivo, de fazer, exprimir e conhecer.
A educação torna -se mais verdadeira, conforme propõe Freire ( 1977 ), à medida que estimula a expressividade dos seres humanos, a Arte é um instrumento de apropriação da realidade descobrindo mundos e recriando -o poeticamente.
O processo de socialização dá -se pelo contato com o outro e, também, pelo contato com a produção do outro ( texto, pintura, música, dança etc. ).
A cultura geral aproxima os homens, pois permite a identificação de uns com os outros.
Para Wallon ( 2008 ), o meio social e a cultura constituem as condições, as possibilidades e os limites do desenvolvimento do organismo.
Os três campos funcionais – o motor, o afetivo e cognitivo são tratados pelo autor de forma indissociável, uma vez que o desenvolvimento de um, necessariamente, causa impactos qualitativos nos demais.
Por isso, o sujeito aprendente precisa ser entendido em seu contexto, e seu desenvolvimento como resultado de sua interação com esse meio.
É importante considerar que as práticas pedagógicas sofrem constantes modificações em busca de novos caminhos.
As metodologias de aprendizado, no campo da Arte/Educação, caminham para novas configurações e propostas que apresentam uma complexidade multissensorial incorporadas e integradas às metodologias ativas e híbridas.
Na metodologia ativa, o estudante é o protagonista e maior responsável pelo processo de aprendizado.
O objetivo desse modelo de ensino é incentivar que os estudantes desenvolvam a capacidade de absorção de conteúdos de maneira autônoma e participativa.
A sala de aula invertida, que tem por objetivo otimizar o tempo em sala, faz com que o estudante chegue com um conhecimento prévio, tire dúvidas com os professores e interaja com os colegas para fazer projetos, resolver problemas ou analisar estudos de caso considerados também métodos ativos.
A metodologia ativa não é a atividade, mas a construção didática que o professor faz e que vai colocar o estudante, cognitivamente e socioemocionalmente, engajado no processo de aprendizagem.
O Ensino Híbrido é uma proposta de educação que se caracteriza por mesclar dois modos de ensino : o on-line e o off-line, momento em que o estudante estuda sozinho, em grupo, com o professor ou colegas, valorizando a interação e o aprendizado coletivo e colaborativo.
Outra metodologia ativa é a \textbf{Espiral} da Aprendizagem, que utiliza tecnologias educacionais, como um recurso para a identificação de problemas que surgem, a partir dos conhecimentos, sentimentos e valores prévios de cada sujeito aprendente.
Os problemas identificados podem ser agrupados, por afinidade, e representam o ponto de partida do processo ensino-aprendizagem.
Não se propõe, aqui, uma receita, mas possibilidades de combinações metodológicas flexíveis que se adequem à realidade de cada contexto escolar.
Pode -se considerar um mundo conectado e digital, com metodologias ativas que se expressam, por meio de modelos de ensinos híbridos, porém o que irá determinar, qual metodologia ou abordagem serão utilizadas, são as especificidades de cada território, no qual a instituição escolar está inserida.
Trabalhar com modelos flexíveis, que utilizem projetos reais, jogos e informação contextualizada, equilibrando a colaboração com a personalização, é o caminho mais significativo, podendo ser planejado e desenvolvido de várias formas e em contextos diferentes.
Cabe ao professor, atuar como condutor e mediador de práticas pedagógicas que potencializem a aprendizagem na sala de aula, de modo que os estudantes possam decodificar e compreender aspectos da cultura de maneira crítica e aprofundada ( AROUCA, 2012, p.9 ).

% matched lemmas: espiral, triangular
\end{textsample}
